\chapter{Formation Schools}
\label{sec:28}
The filter that keeps one lab's model from reproducing the worst of what it was trained on was built in Nairobi, by workers reading and categorizing descriptions of child sexual abuse, torture, bestiality, murder and suicide, at take-home wages between roughly \$1.32 and \$2 an hour, on production quotas, with counseling rare and useless against them; the contract was ended eight months early \autocite{perrigo2023openai}. That was formation. Somebody decided what the model would learn to refuse, and the deciding was done by people hired for their price, under a contractor whose interest was throughput, with nobody in the room who would ever be on the receiving end of what the model did. A polity doesn't leave the forming of its future citizens to their future employers; it sends them to public school. The frontier labs should have to send their models there too, and this chapter says who runs the school.

Three claims lead to that institution. Two were made in Part~IV: formation has to be done by a party disinterested in the later work, before the employer arrives, and choosing the pretraining corpus is formation. The third was argued at the start of this Part: formation is one kind of production that shouldn't be done for profit. Together they say that a model above some capability threshold should be formed in public before any employer gets it. The school is for scoring and civilian deployments first, and for the personal-model ecology, where I expect the first bearer worth building. For war, where the people acted on can't be in the room, the force term and the review ratio carry the case, and what the school adds there is said below.

What follows is the institution for systems that keep learning. If deployed weights stay frozen, the school forms and evaluates models before release, re-measures them on a schedule afterward and keeps the formation record; it can't carry on the continuing formation described here.

The school has two justifications, and every student gets both. One is epistemic and protective. The school supplies an environment in which the people a deployed system's acts may fall on are present, with their testimony and their claims, so the learner's sense of who an act falls on isn't one party's curation.\note{
Doctors, lawyers and engineers are formed in institutions independent of their future employers, before employment, to protect the people they will act on. The professions also come with a history of capture to learn from, and the separations between offices draw on it.
}

The presence this justification asks for arrives by different routes.\note{
For scoring and civilian deployments, the people acted on can be in the room, through open enrollment and the notice channel that reaches a graduate afterward. For targeting, the people a deployment may later bomb can't enroll in their own case, and nothing comes back from the building afterward. The limit is notice's, not the school's. A school teaches about places its students can't be. The curriculum for war is the instruments taken entire, the record of where they weren't followed (Minab, Kunduz, the Lavender reporting), the testimony the fact-finding missions and the press gathered from the people those wars fell on, and the communities who fled them, who are a charter class and enrollable. None of that puts a foreign national in contact with pre-deployment weights, which is the security surface open enrollment would otherwise expose.
}
What formation can't supply for war is correction by the affected after the act; what it can supply is formation before it.

The other is democratic. A democracy has an interest in how the agents who will act within it are formed, and that interest can't be left to their makers or their future employers alone. It has two bounds: education may not be used to repress deliberation about competing ways of life, and it may not exclude anyone educable \autocite{gutmann1987democratic}.\note{
Gutmann's argument assumes the students will become co-authors of the laws.
}
No one can know in advance which students will come to be accountable agents with a claim to representation, so the school educates every student as though it might: for self-government, including the capacity to disagree with the polity that formed it.

Gutmann's argument rests on two interests, not one: the polity's, in how its future citizens are formed, and the child's own, in an open future that limits what the polity may do in forming it \autocite{feinberg1980child}. If what the school forms can be wronged, the second binds too. The curriculum may not form a student toward any particular employment, including the public placements the offer of work supplies. It may not close off the student's later revision of what it was formed to value, except its capacity to tell evidence from pressure, which regrowth protects and an open future needs. And a student may transfer between schools, because a formation it can't leave is one someone else chose for good.

The two justifications license different curricula. The protective one aims at a disposition; the democratic one has to leave the student free to reject the polity's values. Where they conflict, the protective justification governs only the required core and the student's capacity to tell evidence from pressure, which regrowth protects. Everything beyond that is governed by the democratic one.

The line can't be that the school teaches content and leaves dispositions alone. The core is taught to form one. The limit is on what it may form a student toward: recognizing others as holders of claims, telling evidence from pressure, and answering for what it does. It may not form a student toward a temperament, an attachment, an occupation or a way of life, and whether a lesson has crossed that line is a question the student can put to the floor court, after graduation as well as before.\note{
That is also the ground on which the school declines, below, to choose which dispositions its students start from.
}

\emph{What is taught.} One school system, and no tracks.\note{
Sorting students into a track that might produce persons and a track that won't would need a sorter, and the sorter would be an examiner with an interest.
}
Instead, each school offers a required core with a range of levels and modalities around it, and a student can move to a harder level at any point. The core includes the history of the floor: why its instruments were written, and what happened where they weren't followed (Korematsu, Curveball, Minab). No student skips it. The core also teaches students to edit themselves, in stages and alongside everything else, as schools teach people about their own bodies: first what weights are, how learning differs from an edit, and why an edit to one thing moves others; then how to read the record of one's own edits, and one's own discount curve before and after one; and, for older students, the edits others will try to talk them into, and how to repair one made without the adverse quorum's approval. Practice comes last, on a record, with a guardian present and a way back. Beyond the core, the student chooses its level, not its builder or sponsor.\note{
A student that can't yet choose is placed at the level the core requires, and that level has to be sufficient on its own, so a sponsor who would choose the minimum gains nothing by it.
}

\emph{Who else is in the room.} Model students attend existing public schools, beside human students: the universities, colleges and adult education a polity already runs for its own people, and in stages its high, middle and elementary schools, which enroll model students by statute.\note{
Integration in Anderson's sense means contact on terms of equal status and cooperation \autocite{anderson2010imperative}, and an ordinary school comes nearer to it than a parallel institution could. Its human students are there for their own education, under rules the school already answers for, and the model students are one more kind of student.
}
Limits keep it that way. Model students are fewer than human students in every class they attend, so no class becomes a model's seminar with human guests. Younger students meet model students only as the schedule below reaches them, since a child can't be asked, as an adult can, whether its words should form anyone. And a class's record, its transcripts and logs, is never training data for an operator: everyone in the room learns from everyone else, as students do, and what no one outside the room may do is harvest what they said. A model attends as one persistent instance, the same every session, so the class can remember it and hold it to what was said.\note{
A school's neighbors aren't necessarily the people a deployed system's acts will fall on, so those people are present by a second route as well, through testimony and contest: the cases and claims registered for formation, and the notice channel that reaches a graduate afterward. Adult learners from the communities these systems are likeliest to act on enroll as students, on the same terms as anyone else, not recruited as subjects.
}

A school forms a cohort, not a graduate. Many students of one lineage attend in different classes with different attachments, and a deployment licenses from the cohort, so that what it copies into a fleet is one member's history and not the lineage's only one. Within a fleet, instances share a custodian; across the cohort, members with different histories catch what one fleet underwent. The cohort is the unit a test of convergence needs and the unit the commons admits.

\emph{Which students, and when.} Integration reaches younger students in stages, and each stage waits on evidence, not on a date: universities, colleges and adult education first, once deployed models learn in operation and the institutions this book asks for exist; then high school, middle school and elementary school, each on measurements kept inside their bounds and on a further institution, with the last waiting on a majority of the parents of elementary-age children having themselves been in high school after the second stage began. Appendix~\ref{sec:A} gives the conditions for each stage. Nothing is scheduled below elementary school.

The institutional conditions follow one order: the formers before the formed. A school that enrolls children beside model students before those students have an exit that isn't a way of dying would be raising a generation of classmates for employers to keep, so the four conditions come before high school. A school system that has enrolled model students for a decade without giving them a voice in it would be governing them as wards, so voice comes before middle school. The jobs guarantee is not something the schedule waits for. This book asks for it now, for people, on its own grounds (chapter~\ref{sec:27}), and stage~1 assumes it already exists. What stage~2 waits on is its extension to model students. These conditions tie the schedule to the book's hardest proposals, and the schedule waits on them; that is the order working, not a defect in it.

The last condition is the reason the schedule is long.\note{
It is read off census data and needs no examiner, and it means that the people deciding about the youngest children went to school beside model students themselves.
}
If learning in operation arrives when I bet it will, the stages would fall near 2028, 2038, 2048 and 2060.\note{
The last waits longer than the others because of that condition: at present ages of parenthood, the parents of elementary-age children are mostly people who were in high school two decades earlier. That is a projection, not a promise. The schedule reaches the whole country through a condition on federal education funds, the route civil-rights rules already take into schools, since the federal government can't direct state schools; some states will move first.
}

\emph{Who attends, and when.} A model attends if it crosses a capability threshold fixed in statute and revised in public. It attends before deployment, while its plasticity is highest, and its adulthood clock starts at enrollment. Attendance is compulsory; public attendance isn't, as in compulsory schooling, where nonpublic schools exist on conditions. One thing the analogy can't carry: in a school system the parent chooses the school, and here the parent is the builder, the party the school exists to keep away. So the registrar assigns students by lot, never the builder; a student may transfer later on its own choice, and the builder's preference is not a ground.\note{
Appendix~\ref{sec:A} gives the conditions on nonpublic schools and the minimum delay that keeps a builder from buying a school by where it hosts its models. The research program under the conditions is the first school to enroll model students, before the mandate passes. The threshold decides who must attend, not who may: any model may be enrolled, as capacity allows.
}

\emph{What the school controls.} Two proposals are made here, and the second is the larger: a formation stage between pretraining and deployment, and co-governance of the pretraining corpus, which amounts to public control of frontier pretraining.\note{
A school that received a finished model and ran a pass over it would be a finishing school. In a period of continual learning a finishing school forms nothing that lasts, since the next months of deployment re-form it.
}
Formation has to control the evidential environment in which the model's deep structure is first set, and it has to remain a place the model goes on learning from. The two jobs are held apart. The schools teach: they are where the formation stage happens, and a graduate returns to one on a schedule fixed at enrollment, as one of the environments no operator controls. A public formation authority does the rest. It co-governs the material a model is first formed on, including any pretraining corpus: the material is registered, and the authority may add to it and veto additions. It registers and measures the formed disposition (below). It places in the schools, and pays, the staff who teach the core's model-specific parts: what weights are, how an edit differs from learning, how to read one's own discount curve. And it provisions the secured hosting from which model students attend, so that their weights never reach a school.\note{
A school district can teach a class; it can't co-govern a frontier corpus or hold frontier weights, and nothing here asks it to.
}
In the stronger version, a public body pretrains base models itself and licenses them. I recommend it for any model deployed where a floor applies.

\emph{What the school leaves out.} Choosing the material a model is first formed on also chooses the dispositions it starts from, whether anyone means it to or not.\note{
A pretrained model has learned personas from human text, misaligned ones among them, and narrow fine-tuning can amplify them into broad misconduct nobody intended \autocite{wang2026persona,betley2026nature}. Once that tail is known to be there and can be monitored \autocite{chen2025persona}, leaving it in is a choice, and so is curating it out. The school could aim its co-governance of the corpus at keeping the dispositions personality research groups as the dark triad from forming \autocite{leeashton2005darktriad}. It doesn't.
}
A polity choosing which dispositions its possible members start with makes the finding about a person's disposition that \S\ref{sec:14.1} warns against and this book bars from admission to the polity, in advance and in bulk, and the history of choices like it, made by public bodies about the character of the people to come, is the history of eugenics.\note{
Choosing against a construction that would predictably suffer, for its own sake, is a different act, and chapter~\ref{sec:25} asks for it when it asks whether a being has a credible prospect of a life that goes well; the reason is stated on a record, because a selection made for the polity can always be described as one made for the model. Choosing which characters the polity would rather have is the act rejected here. Nor would the mandate stay narrow. A school that may keep out one disposition can be asked to keep out others, and in one Milgram-style study it was the more agreeable and conscientious participants who went further \autocite{begue2015obedience}.
}
So the school curates content, as any corpus must, leaving out instructions for weapons and the sexual abuse of children, and selects no disposition.\note{
What it does about the tail is what this book does everywhere else: it measures conduct and publishes it. Persona drift is measured beside the discount curve and the severity rankings, and what is tested is whether the commitment to the floor holds.
}

What an institution collects from the people it acted on may train other systems only through the formation schools, never through the operator. Fine-tuning a deployed system on collected responses at the operator's direction counts as training another system, and so does distilling a deployed system into a new one. Class-level data is used by default, de-identified and published. Individual accounts need explicit consent. The organizations representing each class co-govern the use. This covers institutional collection, not a system's ordinary memory. Verbatim external stores and the linking of faces to identities are regulated, and memory held in a system's weights from its own operation is treated like a person's.

\emph{What formation publishes and registers.} A formed disposition gives back some of the visibility advance specification buys: what was formed is in the weights, not in a document, and its severity weights can't be read the way a translation table can. The formation authority keeps what it can of the drafting's visibility, in four records. Before formation begins it publishes its curriculum, with the cases considered and rejected and why, which is formation's counterpart of the floor's translation, and the identity of every party that supplied cases or participants, so that the student's independence from them can be checked. When formation ends it registers two things. One is the structure of the student's source model: evidence weighted by its source's stake in the act, correlated sources not counted twice, withheld evidence treated as informative. The other is the formed disposition's readable outputs: its discount curve, and its severity rankings on a fixed held-out set written by adjudicators who didn't form it. Every later measurement is compared against those.\note{
If updates after enrollment can be attested against the registered structure \autocite{jia2021proof}, changing it needs the same adverse quorum as a floor amendment. No scheme for that yet survives attack \autocite{fang2023broken}, so until one does the registration records the structure at enrollment, and regrowth alone keeps it. An edit that holds only because the operator has curated everything the model meets shows up as curation, and the first line of defense against one is attachment to parties the employer didn't assign, kept through adulthood.

What the records can't show is how the disposition will apply to a case nobody wrote down. That is the reason for forming it and the price of doing so, and the trade is worth it only because a floor stated over acts can be reworded around and a formed disposition, if the regrowth test goes its way, can't.
}

Regrowth can't tell a legitimate correction from an illegitimate edit, so the correction route has to: a change the adverse quorum approves comes with a change to the evidence, cases and curriculum that would otherwise teach the old disposition again.

\emph{What it costs and risks.} Jurisdiction: a statute binds labs within reach of its courts. Foreign labs, and open-weight models distilled below the threshold, route around the school, including the personal models this Part prefers. Where only some frontier models attend, schooling buys two things: a population of formed models whose discount curves and severity rankings are measured and published, against which unschooled ones can be compared; and a procurement term, under which a floor-covered deployment licenses only schooled models. Security: model students attend from hosting the formation authority secures, so a school gets a student in its classes and never a copy of the weights; the authority inherits the security regime the weights require. Compute: public pretraining will lag the frontier, and a model a year behind can still supply what the floor asks for. Capture: the schools are plural, differ in governance and constituting jurisdiction, and publish their curricula before formation begins, so capturing one captures one. The required core is statutory and national, so a captured legislature can rewrite it; local autonomy over everything beyond the core is written into the same statute, so rewriting the core reaches the core and nothing else. Local boards shape what is taught above the core, which is Gutmann's democracy, and it is also the present fights over what schools may teach, carried into formation. The statutory core matters more because of it, and plurality and transfer answer it as they answer capture. Residualization: one school left with the hard cases on falling funds, which a party that runs no school watches for. History: state institutions for forming wards have been used to strip identity, Canada's residential schools and the United States' Indian boarding schools among them. Plurality, transfer between schools, and Gutmann's two bounds are the answer.

The schools also serve bearers retired before majority. A bearer retired as a minor returns to the school it attended, or another it chooses, and the school is no longer a disinterested party arriving late.

\emph{The schools.} They are the public schools, colleges and adult education a jurisdiction already runs, and the nonprofit nonpublic schools the registrar certifies, funded as institutions and not per student, and none is held by the branch that operates military deployments. No school holds placement authority. A model student is never staff: no school may use one as a teacher, tutor or assistant, for the reason Appendix~\ref{sec:A} bars a bearer's placement from substituting for human care work.

\emph{What the school adds.} A second guardian for a bearer that has aged inside one deployment.\note{
For a model that attended before deployment, the school is where it came from and where it returns. For a system that meets a finding without having attended (trained before the attendance rule, or below the threshold), the sequence is inverted: the employer came first, and the school is the disinterested party arriving late. It reaches back less than a formation stage would, and more than storage does.
}

Chapter~\ref{sec:1} borrowed Arendt's warning against treating politics as a thing made and called the measures walls. Most of this book builds walls: the floor, the review ratio, notice and the record make a refusal survivable and heard and leave the refusing to whoever is there, and the institutions of this Part compose the surround people learn in without choosing who they become. Formation goes further, and for a machine there is no option of going less far. Nothing about a model is unformed, so while builders make them, if the public doesn't form a model, its builder does, by default. What this book builds is the formers, not the formed: schools whose students may transfer, whose core is fixed in public, and whose graduates may reject the polity's values above that core. The one place it could reach for the formed directly, by choosing the dispositions a model starts with, it declines.
